\documentclass[10pt,a4paper]{amsart}
\usepackage[margin=19mm]{geometry}
\usepackage{amsmath,amssymb,mathtools}
\usepackage[scr=rsfs,cal=euler]{mathalfa}
\usepackage{xfrac}
\usepackage[unicode,hidelinks,bookmarks=false]{hyperref}
\newcommand{\ie}{i.e.\ }
\newcommand{\cf}{cf.\ }
\newcommand{\resp}{respectively\ }
\newcommand{\Subsection}{\subsection}
\providecommand{\nobreakdash}{\nobreak\hbox{-}}
\setlength{\emergencystretch}{2em}
\multlinegap0pt
\usepackage{bm}
\usepackage{bbm}
\usepackage{dsfont}
\usepackage{xspace}
\usepackage{cancel}
\usepackage{mathtools}
\usepackage{amsthm}
\makeatletter
\@ifundefined{theorem}{\newtheorem{theorem}{Theorem}}{}
\@ifundefined{lemma}{\newtheorem{lemma}[theorem]{Lemma}}{}
\makeatother
\title[On the discriminant locus of a generic projection]{On the discriminant locus of a generic projection}
\author{Si-Yang Liu, Yilong Zhang}
\date{Source: arXiv:2603.21518v2 (CC BY 4.0)}
\begin{document}
\maketitle

\noindent\emph{Source and licence.} On the discriminant locus of a generic projection. Version arXiv:2603.21518v2 (CC BY 4.0), distributed under the Creative Commons Attribution 4.0 International licence, as recorded in the arXiv licence metadata for this exact version and shown on the abstract page https://arxiv.org/abs/2603.21518v2. Full rights notice: licenses/arxiv-2603.21518-CC-BY-4.0.txt.\par\smallskip

\noindent\textit{Editorial scope.} Counted unit: the Lemma whose source label is \texttt{lemma\_generalityV} in the article (source0.tex lines 959--994, source label \texttt{lemma\_generalityV}), delivered together with the article's own complete original proof of it (source0.tex lines 996--1036, one \texttt{proof} environment of 41 lines). No mathematics is written, repaired or re-derived: statement and proof are the article's own text line for line. Required context retained uncounted: none. Every definition, display and label of those lines is retained unchanged. Omitted from this package and not redistributed: 1--958, 995--995, 1037--2039. Nothing inside a retained excerpt is omitted or altered: every retained body is delivered exactly as the article's lines are written, so no omission receipt of any kind accompanies this package. Citations: the retained excerpts carry no citation command at all, so no bibliography entry of the article is transcribed and no reference is redistributed. Reference adaptation: none was needed and none was made; every reference inside the delivered unit resolves inside it, so no unresolved reference or citation is produced. Printed slips kept literal: no printed word, symbol, hypothesis or number of the counted statement or of its proof was altered. Layout and class: the article's own class is \texttt{amsart} with packages including inputenc, setspace, babel, mathabx, framed, soul, xy, geometry, arydshln, todonotes, pb-diagram, euscript, multicol, stmaryrd; the wrapper is the lane's pinned \texttt{amsart} editorial wrapper at 10pt on A4, which restates the theorem-like environments the retained text needs and adds the article's own macro definitions that the retained text needs (none), with extra wrapper packages (bm, bbm, dsfont, xspace, cancel, mathtools). The delivered numbering is the unit's own. Assets: no retained span contains a figure environment, a graphics-inclusion command, a PDF-page inclusion, a table or a TikZ picture; no image file is copied into this package, the package declares no asset dependency and no image file is redistributed with it. Licence: version arXiv:2603.21518v2 of this article is distributed under the Creative Commons Attribution 4.0 International licence, as recorded in the arXiv licence metadata for this exact version and shown on the abstract page https://arxiv.org/abs/2603.21518v2. A full rights notice is delivered at licenses/arxiv-2603.21518-CC-BY-4.0.txt. Characters: every retained excerpt is pure ASCII, so the delivered engine, XeLaTeX with its default Latin Modern text font, reports no missing character.
\begin{lemma}[Generality of the linear section]\label{lemma_generalityV}
Let
\[
U:=(N_{\mathbb P^N}^{\perp}X)^\circ
\subset N_{\mathbb P^N}^{\perp}X
\]
be the dense open subset consisting of pairs $(x,\alpha)$ such that
$x\in X^{sm}$, $\alpha\in (X^\perp)^{sm}$, and $T_xX\subseteq \alpha^\perp$.
Then for a general linear subspace $\mathbb PV\subset \mathbb P^N$, the following hold:

\begin{itemize}
\item[(1)] The intersection
$U\cap \pi_1^{-1}(X_V)$
is dense in $\pi_1^{-1}(X_V)$.

\item[(2)] The intersection
$X^{sm}\cap X_V$
is dense in $X_V$.
\end{itemize}

Consequently:

\begin{itemize}
\item[(i)] a general point $\alpha\in (X|_V)^\perp$ satisfies
\[
\alpha\in (X^\perp)^{sm},
\qquad
\exists\, x\in X^{sm}\cap \mathbb PV
\text{ such that } T_xX\subseteq \alpha^\perp;
\]

\item[(ii)] a general point $\alpha_V\in X_V^\perp$ is represented by a hyperplane
$H_V\subset \mathbb PV$ tangent to $X_V$ at some smooth point 
$x\in X_V^{sm}\cap X^{sm}$.
\end{itemize}
\end{lemma}

\begin{proof}
Let $r=N-k$. Choose general hyperplanes $H_1,\dots,H_r\subset \mathbb P^N$
cutting out $\mathbb PV$, and write
$D_i:=\pi_1^*H_i \in |\pi_1^*\mathcal O_{\mathbb P^N}(1)|$. Then
\[
\pi_1^{-1}(X_V)=N_{\mathbb P^N}^{\perp}X\cap D_1\cap\cdots\cap D_r.
\]
Since $U$ is a dense open subset of the irreducible variety
$N_{\mathbb P^N}^{\perp}X$, we may choose the divisors $D_i$ inductively so that
no irreducible component of the intermediate intersections is contained in the
closed subset $N_{\mathbb P^N}^{\perp}X\setminus U$.
It follows that
$U\cap \pi_1^{-1}(X_V)$
is dense in $\pi_1^{-1}(X_V)$, proving (1).
Applying $\pi_2$, whose image is $(X|_V)^\perp$, gives the consequence (i).

Similarly, viewing $X_V=X\cap H_1\cap\cdots\cap H_r$, we may choose the same
general hyperplanes inductively so that no irreducible component of the
intermediate intersections is contained in the closed subset $Sing(X)\subset X$.
Therefore $X^{sm}\cap X_V$
is dense in $X_V$, proving (2).
%Now (i) follows from (1) by applying $\pi_2$:
%a general point of
%\[
%(X|_V)^\perp=\pi_2\big(\pi_1^{-1}(X_V)\big)
%\]
%comes from a pair $(x,\alpha)\in U\cap \pi_1^{-1}(X_V)$.
For (ii), note that
$X_V^{sm}\cap X^{sm}$
is a dense open subset of $X_V$ by (2), hence also dense in $X_V^{sm}$. The rest follows from definition of the dual variety $X_{V}^{\perp}$.
%Therefore the conormal variety
%\[
%N^{\perp}_{\mathbb PV}(X_V)
%=
%\overline{\{(x,H_V)\mid x\in X_V^{sm},\ T_x(X_V)\subseteq H_V\}}
%\]
%is already the closure of the same incidence relation over the dense open subset
%$X_V^{sm}\cap X^{sm}$. Its image in $(\mathbb PV)^\vee$ is $X_V^\perp$, so a
%general point of $X_V^\perp$ is represented by a hyperplane tangent to $X_V$ at
%some point $x\in X_V^{sm}\cap X^{sm}$.
\end{proof}

\end{document}
